\documentclass[10pt,twocolumn]{article}
\usepackage[letterpaper,margin=0.75in,columnsep=0.28in]{geometry}
\usepackage[T1]{fontenc}
\usepackage{lmodern}
\usepackage{microtype}
\usepackage{amsmath,amssymb}
\usepackage{booktabs}
\usepackage{tikz}
\usepackage[font=small,labelfont=bf,skip=4pt]{caption}
\usepackage[hidelinks]{hyperref}
\usepackage[numbers,sort&compress]{natbib}
\setlength{\bibsep}{1pt}
\setlength{\textfloatsep}{8pt plus 2pt minus 2pt}
\setlength{\floatsep}{6pt plus 2pt minus 2pt}
\setlength{\dbltextfloatsep}{8pt plus 2pt minus 2pt}

\newcommand{\VIX}{\mathrm{VIX}}
\newcommand{\E}{\mathbb{E}}
\newcommand{\Q}{\mathbb{Q}}

\title{\vspace{-1.8em}Mean-Reverting Logarithmic Models of the VIX, Revisited:\\
\large A Reproduction of Bao (2013) and an Out-of-Sample Test on the 2026 VIX Option Chain}
\author{Julian Philip Henry}
\date{\vspace{-0.6em}September 2026\vspace{-1.4em}}

\begin{document}
\maketitle

\begin{abstract}
\noindent
Bao (2013) models the logarithm of the VIX as a mean-reverting process, with and without
upward jumps (MRLRJ) and stochastic vol-of-vol (MRLRSV), and calibrates the models to
VIX options from September 2011. Neither the quotes nor the calibration protocol were
published. We give an open implementation, verify it against the thesis'
published parameter tables, and re-run the thesis' calibration procedure unchanged on
the Cboe VIX option chain of 29 September 2026. The central qualitative claims survive:
the pure-diffusion MRLR model cannot produce the VIX skew (pricing errors 28--38\%),
while jumps \emph{or} stochastic vol-of-vol each fit it (errors 1.6--2.8\%, 92--100\%
of model prices inside the bid--ask spread, and similar errors on held-out strikes).
Three further findings qualify them. First, the fitted parameters are largely not
identified: in MRLR the mean-reversion speed has no effect on the fit at all, and the
jump model's per-maturity fits drift to $\kappa\approx0.01$, $\theta\approx-20$.
Second, the 2026 smile needs rare, large jumps ($\lambda\approx2$--$5$/yr, mean log-jump
0.3) or vol-of-vol three to seven times the 2011 estimates, far from the 2011
regime of frequent small jumps. Third, no model fits the four maturities
with one set of parameters: the best (MRLRJ) prices only half the quotes inside the
spread. The models work as per-maturity smile interpolators, not as a
consistent description of VIX dynamics. Code, data snapshot and all fitted parameters
are released.
\end{abstract}

\section{Introduction}
Options on the VIX are among the most liquid volatility derivatives, yet their
underlying is unusual: the index is not traded, it reverts quickly toward a long-run
level, and it jumps up in market stress far more readily than it falls. VIX options
therefore trade with a \emph{positive} implied-volatility skew, the reverse of equity
index options, and VIX futures carry a term structure that is usually in contango and
flips to backwardation in crises.

Bao~\citep{bao2013} builds on the direct-VIX approach of
Whaley~\citep{whaley1993}, Gr\"unbichler and Longstaff~\citep{grunbichler1996} and
Psychoyios et al.~\citep{psychoyios2010}: instead of deriving the VIX from a model of
the S\&P~500, it specifies dynamics for $\ln\VIX_t$ directly and prices futures and
options in closed form or through the characteristic function. The thesis
studies four nested specifications, derives futures, options, hedging ratios and
convexity adjustments for each, and ends with an empirical comparison on Cboe quotes
from 26~September~2011 (spot VIX 42.3). Its conclusions are that the log-normal MRLR
model cannot generate the skew, that adding jumps \emph{or} stochastic vol-of-vol is
enough, and that combining both adds nothing.

The empirical chapter cannot be reproduced as published: the option quotes, the
optimizer, the starting values and the parameter bounds are not reported. This note
(i)~releases an open Julia implementation of MRLR, MRLRJ and MRLRSV; (ii)~tests it
against the parameter tables published in the thesis; (iii)~applies the thesis'
calibration procedure, unchanged, to the Cboe VIX chain of 29~September~2026 (spot
VIX 16.04); and (iv)~asks four questions the original study did not: do the fits
hold on held-out strikes, are they consistent with VIX futures, are the parameters
identified, and can one set of parameters explain the whole term structure?

\section{The models and why each exists}
\label{sec:models}

\paragraph{Why logarithms, why mean reversion.}
A model of the VIX must keep it positive, pull it back toward a long-run level after
spikes, and let its volatility scale with its level. The mean-reverting square-root
(MRSR) process of Gr\"unbichler and Longstaff~\citep{grunbichler1996} satisfies the
first two, but its volatility grows only with $\sqrt{\VIX}$ and it has a thin right
tail. Modelling $X_t=\ln\VIX_t$ as an Ornstein--Uhlenbeck process instead makes the
VIX log-normal with level-proportional volatility:
\begin{equation}
dX_t=\kappa(\theta-X_t)\,dt+\sigma\,dW_t.\tag{MRLR}
\end{equation}
Here $\kappa$ is the speed of mean reversion, $e^\theta$ the long-run VIX level and
$\sigma$ the vol-of-vol. With $\phi=e^{-\kappa\tau}$, $X_T$ is Gaussian with mean
$\mu=\phi X_t+(1-\phi)\theta$ and variance $v=\sigma^2(1-\phi^2)/(2\kappa)$, so the
VIX future is $F=e^{\mu+v/2}$ and a call is a Black formula on $F$. The thesis lets
$\theta_t$ and $\sigma_t$ depend on time so that the model reproduces the futures
curve and the at-the-money volatility term structure by construction.

\paragraph{Why jumps.}
Because MRLR is log-normal, its implied volatility against the future is the same at
every strike: it cannot produce a skew. The VIX's defining feature is sudden upward
jumps, so the thesis (following Psychoyios et al.~\citep{psychoyios2010}) adds
compound-Poisson jumps with exponential sizes:
\begin{equation}
dX_t=\kappa(\theta-X_t)\,dt+\sigma\,dW_t+J\,dN_t,\tag{MRLRJ}
\end{equation}
with $N_t\sim\text{Poisson}(\lambda)$ and $J\sim\text{Exp}(\eta)$, i.e.\ mean
log-jump $1/\eta$. Upward jumps fatten the right tail and create a positive skew,
most strongly at short maturities before mean reversion dampens them. The
characteristic function $\psi(s)=\E^\Q_t[e^{isX_T}]$ stays in closed form:
\begin{equation}
\ln\psi(s)=is\mu-\tfrac12s^2v+\tfrac{\lambda}{\kappa}\ln\frac{\eta-is\phi}{\eta-is}.
\label{eq:cfj}
\end{equation}
The jumps are not compensated, so $F=\psi(-i)$ is finite only if $\eta>1$.

\paragraph{Why stochastic vol-of-vol.}
Vol-of-vol itself varies: it rises when the VIX rises. MRLRSV replaces the constant
$\sigma$ by a square-root variance process correlated with the VIX shocks:
\begin{align}
dX_t&=\kappa(\theta-X_t)\,dt+\sqrt{V_t}\,dW_t,\notag\\
dV_t&=\kappa_v(\theta_v-V_t)\,dt+\sigma_v\sqrt{V_t}\,dZ_t,\tag{MRLRSV}
\end{align}
with $d\langle W,Z\rangle=\rho\,dt$. A positive $\rho$ means that VIX rallies come with
higher vol-of-vol, which fattens the right tail as jumps do but builds up
gradually, so its effect grows with maturity. The model is affine:
$\psi=\exp(a+bV_t+is\phi X_t)$, where, with $u(h)=is\,e^{-\kappa h}$,
\begin{equation}
\begin{aligned}
b'&=\tfrac12u^2+\rho\sigma_vub+\tfrac12\sigma_v^2b^2-\kappa_vb,\\
a&=is\theta(1-\phi)+\kappa_v\theta_v\textstyle\int_0^\tau b\,dh,
\end{aligned}
\label{eq:riccati}
\end{equation}
$b(0)=0$. The Riccati equation has time-dependent coefficients; its Kummer-function
solution is unstable, so the thesis, and we, solve it by Runge--Kutta.

\paragraph{Why not both.}
MRLRSVJ combines jumps and stochastic vol-of-vol. In the thesis it fits no better than
MRLRSV (percentage errors within 0.01 points) at the cost of two more parameters,
so we do not implement it.

\section{Pricing and verification}
Calls follow the Heston-type formula of the thesis (Theorem~4.3), $C=D(F\Pi_1-K\Pi_2)$,
with the Gil-Pelaez~\citep{gilpelaez1951} probabilities
\[
\Pi_2=\tfrac12+\tfrac1\pi\!\int_0^\infty\!\!\frac{\operatorname{Im}\big(e^{-iuk}\psi(u)\big)}{u}du,
\qquad k=\ln K,
\]
and $\Pi_1$ the same with $\psi(u-i)/\psi(-i)$. The integrals use composite 16-point
Gauss--Legendre quadrature on $[0,U]$, with $U$ grown until $|\psi(U)|<10^{-12}$ and
one grid shared by all strikes of a maturity. For MRLRSV, (\ref{eq:riccati}) is solved
by RK4 at each node, with a step that resolves $e^{-\kappa h}$ and keeps
$h(\sigma_v|s|+\kappa_v)\le1$ inside the stability region of this stiff equation.

A test suite (34 assertions) checks that the Fourier pricer reproduces the MRLR closed
form to $10^{-8}$; that MRLRJ ($\lambda\to0$) and MRLRSV ($\sigma_v\to0$,
$V_0=\theta_v$) reduce to MRLR; that $F=\psi(-i)$ matches the closed-form futures;
that prices are decreasing, convex and within arbitrage bounds; that Euler Monte
Carlo prices (200k paths) agree within four standard errors for MRLRJ and MRLRSV;
and that a tenfold smaller RK4 step changes prices by less than $10^{-4}$. A
25-digit \texttt{BigFloat} quadrature agrees with the double-precision pricer to
$10^{-9}$, so extended precision is unnecessary.

\section{Consistency with the published tables}
The 2011 quotes are not public, so the thesis' errors cannot be recomputed. Its
parameter tables can still be tested. Tables~7.5 (MRLRJ) and 7.6 (MRLRSV) were fitted
separately to the \emph{same} quotes, with reported mean absolute errors (MAE)
of 0.12--0.21 each. Priced with our code ($\VIX_t=42.3$, strikes 20--90), the two
structurally different models agree to 2--3.5 cents at every maturity
(Table~\ref{tab:2011}), well inside that tolerance; MRLR (Table~7.4) differs by
8--17 cents. Against the model future, the MRLR implied volatility is flat and the
MRLRJ and MRLRSV curves slope upward, at the levels plotted in the thesis'
Figure~7.3 (e.g.\ 1.27/1.39/1.48 for MRLRJ at $K=40/60/80$ in October). A
convention error in the jump or Riccati terms would break this agreement, so the
implementation matches the thesis.

The tables also imply numbers the thesis never reports: model VIX futures of
37.8--38.5 (Oct-2011), 35.0--35.4 (Nov), 31.7--32.6 (Dec) and 32.2--33.4 (Jan-2012),
a backwardated curve consistent with the text. They also show that the calibration
ran into undisclosed bounds: $\theta=3.00$ in all four MRLRJ fits and $\rho=1.00$ in
all four MRLRSV fits are boundary values, not estimates.

\begin{table}[t]
\centering\small
\caption{Published 2011 parameters priced with our implementation: mean absolute
price difference between models (strikes 20--90) and the thesis' reported MAE
against the market (Table~7.3).}
\label{tab:2011}
\footnotesize\setlength{\tabcolsep}{3pt}
\begin{tabular}{@{}lcccc@{}}
\toprule
Expiry & $|\mathrm{J}-\mathrm{SV}|$ & MAE$_\mathrm{J}{+}$MAE$_\mathrm{SV}$ & $|\mathrm{LR}-\mathrm{J}|$ & MAE$_\mathrm{LR}$\\
\midrule
18 Oct 2011 & 0.035 & 0.238 & 0.173 & 0.433\\
15 Nov 2011 & 0.019 & 0.359 & 0.105 & 0.329\\
20 Dec 2011 & 0.022 & 0.409 & 0.079 & 0.357\\
17 Jan 2012 & 0.022 & 0.261 & 0.121 & 0.357\\
\bottomrule
\end{tabular}
\end{table}

\section{Data and calibration protocol}
\label{sec:protocol}
\paragraph{Data.} We use the Cboe delayed-quote snapshot of the full VIX option chain
at the close of 29~September~2026 (spot VIX 16.04 at 16:15 ET; 1{,}520 contracts;
SHA-256 \texttt{c404ecd5\ldots ae5f9})~\citep{cboe}. As in the thesis we take the first
four monthly maturities (21~Oct, 18~Nov, 16~Dec 2026 and 20~Jan 2027; open interest
falls sharply after January), keep calls with positive bid \emph{and} positive open
interest (55, 63, 58 and 59 quotes) and fit mid quotes. Time to expiry runs to the
09:30~ET settlement on the expiry date (ACT/365).

\paragraph{Futures and rate.} VIX options settle on the same quotation as VIX
futures, so put--call parity holds against the future:
$C-P=e^{-r\tau}(F-K)$. Fitting this jointly over all call--put pairs gives
$r=4.46\%$ and futures of 17.59, 18.28, 18.67 and 19.41, a contango curve
above the spot. The thesis instead used spot VIX with no discounting adjustment.

\paragraph{Loss and fit measures.} Exactly as in the thesis (Eq.~7.3):
$\mathcal L=\sum_j(C^\text{mod}_j-C^\text{mkt}_j)^2+\alpha\sum_j(\ln C^\text{mod}_j-\ln C^\text{mkt}_j)^2$
with $\alpha=8$, which balances in- and out-of-the-money quotes. We report the
percentage error $\text{PE}=\text{mean}\,|C^\text{mod}-C^\text{mkt}|/C^\text{mkt}$, the
MAE and, as a stricter test, the share of model prices inside the bid--ask band.

\paragraph{Optimisation.} This is the part the thesis omits.
Parameters are mapped to unconstrained coordinates: $\log$ for positive parameters,
$\eta=1+e^x$, $\rho=\tanh x$. MRLRSV uses logistic box bounds
$\kappa\in[0.1,50.1]$, $\theta\in[1,5]$, $\kappa_v,\theta_v\in(0,50)$ and
$\sigma_v,V_0\in(0,30)$, which rule out a degenerate ridge ($\kappa\to0$,
$\theta\to-\infty$) and extreme stiffness. Our first bounds, $\sigma_v,V_0\le5$, were
binding on the October slice and were widened, since a 2026 short-dated smile needs
far more vol-of-vol than in 2011. Each fit runs Nelder--Mead from 10 random starts
(6 for MRLRSV), with seeds $1+k$ and log-uniform ranges, followed by one restart from
the optimum; each run is capped at 3{,}000 iterations and 300\,s. The ``paper
procedure'' fits every parameter, including $\theta$, per maturity with spot VIX
as input, as in the thesis.

\section{Results on the 2026 chain}

\paragraph{Fit quality.} Table~\ref{tab:fit} repeats the thesis' Tables~7.2--7.3 on the
2026 data. The ranking is the same and the gaps are much larger. MRLR misses by
28--38\% and prices only about a fifth of the quotes inside the spread. MRLRJ and
MRLRSV are essentially tied, with errors of 1.6--2.8\% (lower than the thesis'
3.0--5.2\% in 2011) and 92--100\% of prices inside the spread. As in 2011, the
difference between jump and stochastic-vol models is within noise.
Figure~\ref{fig:skew} shows why: the 2026 skew is steep, rising from about 0.7 in implied volatility at the money to about 2.0 at $K=50$ for the
October expiry.
MRLR's smile is flat against its \emph{own} future, so it can match either the
level or the slope of this curve but not both.

\begin{table}[t]
\centering\small
\caption{Paper procedure on 29 Sep 2026: percentage error, MAE and share of
model prices inside the bid--ask spread. Thesis (2011) PE: MRLR 4.7--6.9\%,
MRLRJ 3.1--5.2\%, MRLRSV 3.0--5.1\%.}
\label{tab:fit}
\setlength{\tabcolsep}{3pt}
\begin{tabular}{@{}l rrr rrr rrr@{}}
\toprule
 & \multicolumn{3}{c}{PE (\%)} & \multicolumn{3}{c}{MAE} & \multicolumn{3}{c}{in spread (\%)}\\
\cmidrule(lr){2-4}\cmidrule(lr){5-7}\cmidrule(l){8-10}
Expiry & LR & J & SV & LR & J & SV & LR & J & SV\\
\midrule
Oct-26 & 37.9 & 1.9 & 1.6 & .527 & .007 & .008 & 18 & 100 & 100\\
Nov-26 & 29.5 & 2.7 & 2.8 & .517 & .016 & .019 & 21 & 98 & 98\\
Dec-26 & 30.6 & 2.3 & 2.4 & .516 & .014 & .017 & 22 & 98 & 98\\
Jan-27 & 28.3 & 2.4 & 2.4 & .527 & .020 & .021 & 22 & 93 & 92\\
\bottomrule
\end{tabular}
\end{table}

\begin{figure*}[t]
\centering
\begin{tikzpicture}[font=\scriptsize]
\draw[black!40] (0,0) rectangle (6.300,4.200);
\draw (0.000,0) -- (0.000,-0.08) node[below] {10};
\draw (0.787,0) -- (0.787,-0.08) node[below] {15};
\draw (1.575,0) -- (1.575,-0.08) node[below] {20};
\draw (2.362,0) -- (2.362,-0.08) node[below] {25};
\draw (3.150,0) -- (3.150,-0.08) node[below] {30};
\draw (3.938,0) -- (3.938,-0.08) node[below] {35};
\draw (4.725,0) -- (4.725,-0.08) node[below] {40};
\draw (5.513,0) -- (5.513,-0.08) node[below] {45};
\draw (6.300,0) -- (6.300,-0.08) node[below] {50};
\draw (0,0.000) -- (-0.08,0.000) node[left] {0.3};
\draw (0,0.933) -- (-0.08,0.933) node[left] {0.7};
\draw (0,1.867) -- (-0.08,1.867) node[left] {1.1};
\draw (0,2.800) -- (-0.08,2.800) node[left] {1.5};
\draw (0,3.733) -- (-0.08,3.733) node[left] {1.9};
\node at (3.150,-0.55) {strike};
\node[rotate=90] at (-0.75,2.100) {Black implied vol.};
\node[anchor=north west] at (0.05,4.180) {\textbf{21 Oct 2026} ($\tau$=0.060, $F$=17.59)};
\begin{scope}\clip (0,0) rectangle (6.3,4.2);
\draw[black!45, line width=1.6pt] (0.787,0.124) -- (0.787,0.620);
\draw[black!45, line width=1.6pt] (0.866,0.335) -- (0.866,0.609);
\draw[black!45, line width=1.6pt] (0.945,0.485) -- (0.945,0.688);
\draw[black!45, line width=1.6pt] (1.024,0.657) -- (1.024,0.812);
\draw[black!45, line width=1.6pt] (1.102,0.875) -- (1.102,0.975);
\draw[black!45, line width=1.6pt] (1.181,0.994) -- (1.181,1.118);
\draw[black!45, line width=1.6pt] (1.260,1.161) -- (1.260,1.270);
\draw[black!45, line width=1.6pt] (1.339,1.290) -- (1.339,1.400);
\draw[black!45, line width=1.6pt] (1.417,1.458) -- (1.417,1.543);
\draw[black!45, line width=1.6pt] (1.496,1.533) -- (1.496,1.648);
\draw[black!45, line width=1.6pt] (1.575,1.689) -- (1.575,1.762);
\draw[black!45, line width=1.6pt] (1.654,1.754) -- (1.654,1.860);
\draw[black!45, line width=1.6pt] (1.733,1.833) -- (1.733,1.958);
\draw[black!45, line width=1.6pt] (1.811,1.979) -- (1.811,2.043);
\draw[black!45, line width=1.6pt] (1.890,1.999) -- (1.890,2.116);
\draw[black!45, line width=1.6pt] (1.969,2.074) -- (1.969,2.195);
\draw[black!45, line width=1.6pt] (2.047,2.139) -- (2.047,2.264);
\draw[black!45, line width=1.6pt] (2.126,2.194) -- (2.126,2.324);
\draw[black!45, line width=1.6pt] (2.205,2.260) -- (2.205,2.394);
\draw[black!45, line width=1.6pt] (2.284,2.318) -- (2.284,2.457);
\draw[black!45, line width=1.6pt] (2.362,2.389) -- (2.362,2.513);
\draw[black!45, line width=1.6pt] (2.520,2.493) -- (2.520,2.625);
\draw[black!45, line width=1.6pt] (2.677,2.572) -- (2.677,2.737);
\draw[black!45, line width=1.6pt] (2.835,2.678) -- (2.835,2.830);
\draw[black!45, line width=1.6pt] (2.992,2.768) -- (2.992,2.929);
\draw[black!45, line width=1.6pt] (3.150,2.842) -- (3.150,3.014);
\draw[black!45, line width=1.6pt] (3.308,2.931) -- (3.308,3.084);
\draw[black!45, line width=1.6pt] (3.465,3.009) -- (3.465,3.170);
\draw[black!45, line width=1.6pt] (3.622,3.074) -- (3.622,3.278);
\draw[black!45, line width=1.6pt] (3.780,3.128) -- (3.780,3.345);
\draw[black!45, line width=1.6pt] (3.938,3.211) -- (3.938,3.438);
\draw[black!45, line width=1.6pt] (4.095,3.287) -- (4.095,3.486);
\draw[black!45, line width=1.6pt] (4.253,3.356) -- (4.253,3.564);
\draw[black!45, line width=1.6pt] (4.410,3.370) -- (4.410,3.636);
\draw[black!45, line width=1.6pt] (4.567,3.423) -- (4.567,3.702);
\draw[black!45, line width=1.6pt] (4.725,3.469) -- (4.725,3.762);
\draw[black!45, line width=1.6pt] (5.119,3.583) -- (5.119,3.911);
\draw[black!45, line width=1.6pt] (5.513,3.658) -- (5.513,4.030);
\draw[black!45, line width=1.6pt] (5.906,3.775) -- (5.906,4.121);
\draw[black!45, line width=1.6pt] (6.300,3.773) -- (6.300,4.182);
\draw[black!55, densely dotted, thick] (0.866,0.572) -- (0.945,1.019) -- (1.024,1.312) -- (1.102,1.539) -- (1.181,1.726) -- (1.260,1.886) -- (1.339,2.024) -- (1.417,2.147) -- (1.496,2.257) -- (1.575,2.356) -- (1.654,2.446) -- (1.733,2.529) -- (1.811,2.604) -- (1.890,2.675) -- (1.969,2.740) -- (2.047,2.800) -- (2.126,2.857) -- (2.205,2.910) -- (2.284,2.960) -- (2.362,3.006) -- (2.520,3.092) -- (2.677,3.170) -- (2.835,3.239) -- (2.992,3.302) -- (3.150,3.360) -- (3.308,3.413) -- (3.465,3.462) -- (3.622,3.507) -- (3.780,3.548) -- (3.938,3.587) -- (4.095,3.624) -- (4.253,3.658) -- (4.410,3.689) -- (4.567,3.719) -- (4.725,3.747) -- (5.119,3.811) -- (5.513,3.867) -- (5.906,3.917) -- (6.300,3.961);
\draw[red!75!black, solid, thick] (0.472,0.018) -- (0.551,0.212) -- (0.630,0.269) -- (0.709,0.332) -- (0.787,0.408) -- (0.866,0.502) -- (0.945,0.617) -- (1.024,0.751) -- (1.102,0.898) -- (1.181,1.051) -- (1.260,1.200) -- (1.339,1.342) -- (1.417,1.472) -- (1.496,1.591) -- (1.575,1.700) -- (1.654,1.801) -- (1.733,1.894) -- (1.811,1.981) -- (1.890,2.063) -- (1.969,2.140) -- (2.047,2.212) -- (2.126,2.281) -- (2.205,2.347) -- (2.284,2.409) -- (2.362,2.469) -- (2.520,2.581) -- (2.677,2.684) -- (2.835,2.781) -- (2.992,2.870) -- (3.150,2.954) -- (3.308,3.034) -- (3.465,3.109) -- (3.622,3.180) -- (3.780,3.247) -- (3.938,3.311) -- (4.095,3.372) -- (4.253,3.431) -- (4.410,3.487) -- (4.567,3.541) -- (4.725,3.592) -- (5.119,3.713) -- (5.513,3.824) -- (5.906,3.926) -- (6.300,4.020);
\draw[blue!70!black, dashed, thick] (0.472,0.505) -- (0.551,0.532) -- (0.630,0.523) -- (0.709,0.502) -- (0.787,0.482) -- (0.866,0.479) -- (0.945,0.526) -- (1.024,0.688) -- (1.102,0.889) -- (1.181,1.068) -- (1.260,1.226) -- (1.339,1.363) -- (1.417,1.488) -- (1.496,1.604) -- (1.575,1.706) -- (1.654,1.807) -- (1.733,1.896) -- (1.811,1.982) -- (1.890,2.063) -- (1.969,2.138) -- (2.047,2.212) -- (2.126,2.278) -- (2.205,2.346) -- (2.284,2.405) -- (2.362,2.467) -- (2.520,2.579) -- (2.677,2.681) -- (2.835,2.776) -- (2.992,2.866) -- (3.150,2.955) -- (3.308,3.037) -- (3.465,3.108) -- (3.622,3.175) -- (3.780,3.249) -- (3.938,3.316) -- (4.095,3.368) -- (4.253,3.434) -- (4.410,3.492) -- (4.567,3.536) -- (4.725,3.601) -- (5.119,3.707) -- (5.513,3.833) -- (5.906,3.933) -- (6.300,4.011);
\end{scope}
\draw[black!45, line width=1.6pt] (4.400,0.970) -- (4.400,1.150) node[right, black] {bid--ask};
\draw[black!55, densely dotted, thick] (4.200,0.770) -- (4.600,0.770) node[right, black] {MRLR};
\draw[red!75!black, solid, thick] (4.200,0.490) -- (4.600,0.490) node[right, black] {MRLRJ};
\draw[blue!70!black, dashed, thick] (4.200,0.210) -- (4.600,0.210) node[right, black] {MRLRSV};
\end{tikzpicture}
\hspace{0.9cm}
\begin{tikzpicture}[font=\scriptsize]
\draw[black!40] (0,0) rectangle (6.300,4.200);
\draw (0.000,0) -- (0.000,-0.08) node[below] {10};
\draw (1.260,0) -- (1.260,-0.08) node[below] {20};
\draw (2.520,0) -- (2.520,-0.08) node[below] {30};
\draw (3.780,0) -- (3.780,-0.08) node[below] {40};
\draw (5.040,0) -- (5.040,-0.08) node[below] {50};
\draw (6.300,0) -- (6.300,-0.08) node[below] {60};
\draw (0,0.000) -- (-0.08,0.000) node[left] {0.3};
\draw (0,0.840) -- (-0.08,0.840) node[left] {0.5};
\draw (0,1.680) -- (-0.08,1.680) node[left] {0.7};
\draw (0,2.520) -- (-0.08,2.520) node[left] {0.9};
\draw (0,3.360) -- (-0.08,3.360) node[left] {1.1};
\draw (0,4.200) -- (-0.08,4.200) node[left] {1.3};
\node at (3.150,-0.55) {strike};
\node[rotate=90] at (-0.75,2.100) {Black implied vol.};
\node[anchor=north west] at (0.05,4.180) {\textbf{20 Jan 2027} ($\tau$=0.309, $F$=19.41)};
\begin{scope}\clip (0,0) rectangle (6.3,4.2);
\draw[black!45, line width=1.6pt] (0.441,0.094) -- (0.441,0.908);
\draw[black!45, line width=1.6pt] (0.504,0.234) -- (0.504,0.956);
\draw[black!45, line width=1.6pt] (0.567,0.539) -- (0.567,0.962);
\draw[black!45, line width=1.6pt] (0.630,0.577) -- (0.630,0.940);
\draw[black!45, line width=1.6pt] (0.756,0.835) -- (0.756,1.108);
\draw[black!45, line width=1.6pt] (0.882,1.028) -- (0.882,1.262);
\draw[black!45, line width=1.6pt] (1.008,1.261) -- (1.008,1.422);
\draw[black!45, line width=1.6pt] (1.134,1.439) -- (1.134,1.602);
\draw[black!45, line width=1.6pt] (1.260,1.596) -- (1.260,1.726);
\draw[black!45, line width=1.6pt] (1.386,1.749) -- (1.386,1.868);
\draw[black!45, line width=1.6pt] (1.512,1.878) -- (1.512,1.997);
\draw[black!45, line width=1.6pt] (1.638,2.002) -- (1.638,2.112);
\draw[black!45, line width=1.6pt] (1.764,2.106) -- (1.764,2.218);
\draw[black!45, line width=1.6pt] (1.890,2.205) -- (1.890,2.320);
\draw[black!45, line width=1.6pt] (2.016,2.293) -- (2.016,2.410);
\draw[black!45, line width=1.6pt] (2.142,2.384) -- (2.142,2.503);
\draw[black!45, line width=1.6pt] (2.268,2.469) -- (2.268,2.580);
\draw[black!45, line width=1.6pt] (2.394,2.539) -- (2.394,2.653);
\draw[black!45, line width=1.6pt] (2.520,2.607) -- (2.520,2.724);
\draw[black!45, line width=1.6pt] (2.646,2.675) -- (2.646,2.795);
\draw[black!45, line width=1.6pt] (2.772,2.732) -- (2.772,2.855);
\draw[black!45, line width=1.6pt] (2.898,2.791) -- (2.898,2.918);
\draw[black!45, line width=1.6pt] (3.024,2.855) -- (3.024,2.972);
\draw[black!45, line width=1.6pt] (3.150,2.910) -- (3.150,3.030);
\draw[black!45, line width=1.6pt] (3.276,2.958) -- (3.276,3.082);
\draw[black!45, line width=1.6pt] (3.402,3.013) -- (3.402,3.126);
\draw[black!45, line width=1.6pt] (3.528,3.062) -- (3.528,3.177);
\draw[black!45, line width=1.6pt] (3.654,3.104) -- (3.654,3.223);
\draw[black!45, line width=1.6pt] (3.780,3.141) -- (3.780,3.278);
\draw[black!45, line width=1.6pt] (4.095,3.242) -- (4.095,3.372);
\draw[black!45, line width=1.6pt] (4.410,3.348) -- (4.410,3.418);
\draw[black!45, line width=1.6pt] (4.725,3.431) -- (4.725,3.559);
\draw[black!45, line width=1.6pt] (5.040,3.492) -- (5.040,3.647);
\draw[black!45, line width=1.6pt] (5.670,3.618) -- (5.670,3.792);
\draw[black!45, line width=1.6pt] (6.300,3.727) -- (6.300,3.922);
\draw[black!55, densely dotted, thick] (0.630,0.482) -- (0.756,1.193) -- (0.882,1.605) -- (1.008,1.905) -- (1.134,2.141) -- (1.260,2.334) -- (1.386,2.497) -- (1.512,2.637) -- (1.638,2.758) -- (1.764,2.866) -- (1.890,2.962) -- (2.016,3.048) -- (2.142,3.125) -- (2.268,3.196) -- (2.394,3.260) -- (2.520,3.320) -- (2.646,3.374) -- (2.772,3.425) -- (2.898,3.472) -- (3.024,3.516) -- (3.150,3.557) -- (3.276,3.595) -- (3.402,3.631) -- (3.528,3.665) -- (3.654,3.697) -- (3.780,3.727) -- (4.095,3.796) -- (4.410,3.857) -- (4.725,3.911) -- (5.040,3.959) -- (5.670,4.042) -- (6.300,4.111);
\draw[red!75!black, solid, thick] (0.000,0.711) -- (0.063,0.580) -- (0.126,0.466) -- (0.189,0.385) -- (0.252,0.351) -- (0.315,0.364) -- (0.378,0.410) -- (0.441,0.478) -- (0.504,0.561) -- (0.567,0.653) -- (0.630,0.754) -- (0.756,0.968) -- (0.882,1.185) -- (1.008,1.389) -- (1.134,1.572) -- (1.260,1.733) -- (1.386,1.875) -- (1.512,2.001) -- (1.638,2.114) -- (1.764,2.216) -- (1.890,2.310) -- (2.016,2.396) -- (2.142,2.476) -- (2.268,2.551) -- (2.394,2.621) -- (2.520,2.686) -- (2.646,2.748) -- (2.772,2.807) -- (2.898,2.862) -- (3.024,2.915) -- (3.150,2.965) -- (3.276,3.013) -- (3.402,3.059) -- (3.528,3.103) -- (3.654,3.145) -- (3.780,3.186) -- (4.095,3.280) -- (4.410,3.367) -- (4.725,3.447) -- (5.040,3.522) -- (5.670,3.655) -- (6.300,3.773);
\draw[blue!70!black, dashed, thick] (0.189,0.353) -- (0.252,0.444) -- (0.315,0.491) -- (0.378,0.526) -- (0.441,0.558) -- (0.504,0.595) -- (0.567,0.640) -- (0.630,0.697) -- (0.756,0.861) -- (0.882,1.132) -- (1.008,1.378) -- (1.134,1.575) -- (1.260,1.740) -- (1.386,1.882) -- (1.512,2.008) -- (1.638,2.120) -- (1.764,2.222) -- (1.890,2.316) -- (2.016,2.402) -- (2.142,2.482) -- (2.268,2.557) -- (2.394,2.627) -- (2.520,2.692) -- (2.646,2.752) -- (2.772,2.811) -- (2.898,2.868) -- (3.024,2.919) -- (3.150,2.968) -- (3.276,3.018) -- (3.402,3.063) -- (3.528,3.106) -- (3.654,3.150) -- (3.780,3.188) -- (4.095,3.285) -- (4.410,3.369) -- (4.725,3.450) -- (5.040,3.527) -- (5.670,3.658) -- (6.300,3.774);
\end{scope}
\draw[black!45, line width=1.6pt] (4.400,0.970) -- (4.400,1.150) node[right, black] {bid--ask};
\draw[black!55, densely dotted, thick] (4.200,0.770) -- (4.600,0.770) node[right, black] {MRLR};
\draw[red!75!black, solid, thick] (4.200,0.490) -- (4.600,0.490) node[right, black] {MRLRJ};
\draw[blue!70!black, dashed, thick] (4.200,0.210) -- (4.600,0.210) node[right, black] {MRLRSV};
\end{tikzpicture}

\caption{Black implied volatility against the parity-implied VIX future for the
nearest and farthest calibrated maturities, 29~Sep~2026. Grey bars: market bid--ask;
lines: models calibrated per maturity with the thesis' procedure. MRLR is flat
against its own future, which is 27--30\% below the market (Sec.~6), so against the market
future its curve is misplaced; MRLRJ and MRLRSV follow the steep positive skew.}
\label{fig:skew}
\end{figure*}

\paragraph{Held-out strikes.} Refitting on strikes of odd rank and pricing the rest
(Table~\ref{tab:oos}) leaves MRLRJ almost unchanged (1.9--2.6\% except 4.7\% in
January). MRLRSV degrades a little more (2.2--3.5\%). The fits interpolate the
smile; they are not memorising individual quotes.

\begin{table}[t]
\centering\small
\caption{Out-of-sample strikes: fit on odd-ranked strikes, evaluate on the rest.
PE in \%; last columns: share inside the spread.}
\label{tab:oos}
\setlength{\tabcolsep}{4pt}
\begin{tabular}{@{}l rrr rr@{}}
\toprule
 & \multicolumn{3}{c}{out-of-sample PE} & \multicolumn{2}{c}{in spread (\%)}\\
\cmidrule(lr){2-4}\cmidrule(l){5-6}
Expiry & MRLR & MRLRJ & MRLRSV & J & SV\\
\midrule
Oct-26 & 39.6 & 1.9 & 3.5 & 100 & 96\\
Nov-26 & 30.2 & 2.6 & 2.6 & 100 & 94\\
Dec-26 & 31.9 & 2.2 & 2.2 & 100 & 93\\
Jan-27 & 29.3 & 4.7 & 3.2 & 90 & 90\\
\bottomrule
\end{tabular}
\end{table}

\paragraph{Futures consistency.} The thesis' procedure uses spot VIX as the underlying
and never looks at futures, which the thesis itself flags as problematic (Sec.~7.3.2). In 2026
this matters only for MRLR. The MRLRJ and MRLRSV fits reproduce the parity-implied
futures to within 0.1\% although the futures were never targeted. Deep in-the-money
calls pin the forward, and the free per-maturity $\theta$ absorbs the gap between spot
and future. The MRLR fits put the future 27--30\% below the market (12.8 vs.\ 17.6 in
October): to fit the right tail through the log-loss, the log-normal model gives up
the level.

\paragraph{Identification.} The parameters are far less meaningful than the fit quality
suggests. In MRLR, one maturity determines only the mean and variance of $X_T$, so
$(\kappa,\theta,\sigma)$ has one redundant direction. Fixing $\kappa$ at 0.5, 2, 8 or
32 and refitting $\theta,\sigma$ gives the same loss to nine digits (120.910550). The
thesis' MRLR values $\kappa=11.05$, 11.43, 9.88 and 12.58 are therefore arbitrary
points on a flat ridge. The same ridge appears in the 2026 MRLRJ fits: beyond
October the optimum drifts to $\kappa\approx0.01$ and $\theta\approx-15$ to $-23$
(Table~\ref{tab:params}). Only the product $\kappa\theta\approx-0.21$ to $-0.26$ is
pinned down, so the jump model degenerates into a drifting Brownian motion with jumps,
with no detectable mean reversion within any single maturity. In MRLRSV, $\kappa_v$ and
$\theta_v$ collapse toward zero at three maturities, so the variance process
effectively has no mean reversion either, and in 2011 $\theta$ and $\rho$ sat on
bounds. Per-maturity MRLR-family parameters should not be read as
estimates of VIX dynamics.

\paragraph{A different volatility regime.} The parameters that are identified have
moved a long way since 2011. There, MRLRJ used many small jumps
($\lambda=60$--$170$/yr, mean log-jump $1/\eta\approx0.10$--$0.15$), close to extra
diffusion. In 2026 the October fit uses rare large jumps: $\lambda=5.2$/yr, mean log-jump
$0.30$ (a 35\% move in the VIX), about 0.3 jumps expected before expiry, on top of a
small diffusion $\sigma=0.35$. Likewise MRLRSV needs $\sigma_v=3.2$--$14.5$ against
0.57--1.98 in 2011. At a spot of 16, far below the 42 of 2011, the market prices a
calm VIX with a large, discrete crash-spike risk.

\paragraph{One model for the term structure.} Finally, we fit a \emph{single} parameter set
to all four maturities, with $\theta$ set per maturity so that each model future
matches the market exactly. This is the thesis' own two-stage strategy (Theorems
3.3, 4.4, 5.4) with constant vol-of-vol. Table~\ref{tab:joint} shows that the per-maturity
success does not carry over. MRLRJ, the best, misses by 5--13\% and prices only
52\% of quotes inside the spread; MRLRSV does worse (12--31\%); MRLR fails
(69--95\%). The MRLRSV run hit its time limit, so its figures are an upper bound, but
the gap to the per-maturity fits is far larger than any plausible optimizer slack.
The thesis anticipates this by making $\sigma_t$ time-dependent, which effectively
reintroduces per-maturity parameters.

\begin{table}[t]
\centering\small
\caption{One parameter set for all four maturities ($\theta$ anchored to each future).
PE in \%. $^\ast$time-capped optimisation.}
\label{tab:joint}
\setlength{\tabcolsep}{4pt}
\begin{tabular}{@{}l rrrr r@{}}
\toprule
Model & Oct & Nov & Dec & Jan & in spread\\
\midrule
MRLR & 95.0 & 81.0 & 71.5 & 68.7 & 9\%\\
MRLRJ & 12.6 & 6.8 & 5.3 & 6.3 & 52\%\\
MRLRSV$^\ast$ & 16.1 & 12.1 & 29.0 & 31.1 & 35\%\\
\bottomrule
\end{tabular}
\\[3pt]
\footnotesize MRLRJ: $\kappa=1.20$, $\sigma=0.469$, $\lambda=2.33$, $\eta=3.12$;
MRLRSV: $\kappa=1.80$, $\rho=0.96$, $\kappa_v=18.1$, $\theta_v=0.18$, $\sigma_v=8.17$, $V_0=1.59$.
\end{table}

\begin{table*}[t]
\centering\small
\caption{Full disclosure: all parameters from the thesis' per-maturity procedure on
29 Sep 2026 (spot VIX 16.04, $r=4.46\%$). Values at or near a bound are
marked $^\dagger$. Per-strike fits: \texttt{results/calibration\_2026-09-29\_curves.csv}.}
\label{tab:params}
\setlength{\tabcolsep}{5pt}
\begin{tabular}{@{}l ccc ccccc ccccccc@{}}
\toprule
 & \multicolumn{3}{c}{MRLR} & \multicolumn{5}{c}{MRLRJ} & \multicolumn{7}{c}{MRLRSV}\\
\cmidrule(lr){2-4}\cmidrule(lr){5-9}\cmidrule(l){10-16}
Expiry & $\kappa$ & $\theta$ & $\sigma$ & $\kappa$ & $\theta$ & $\sigma$ & $\lambda$ & $\eta$
 & $\kappa$ & $\theta$ & $\rho$ & $\kappa_v$ & $\theta_v$ & $\sigma_v$ & $V_0$\\
\midrule
Oct-26 & 88.5 & 2.35 & 8.35 & 9.29 & 2.76 & 0.353 & 5.23 & 3.29 & 29.1 & 2.86 & 0.93 & 1.31 & 13.6 & 14.5 & 1.74\\
Nov-26 & 19.6 & 2.17 & 5.15 & 0.015 & $-15.0$ & 0.263 & 2.93 & 3.45 & 3.48 & 2.98 & 0.75 & 0.20 & $0^\dagger$ & 5.39 & 0.92\\
Dec-26 & 0.265 & $-8.15$ & 1.99 & 0.014 & $-18.7$ & 0.265 & 2.38 & 3.42 & 2.07 & 3.00 & 0.71 & $0^\dagger$ & $0^\dagger$ & 3.48 & 0.75\\
Jan-27 & 0.325 & $-2.68$ & 1.69 & 0.009 & $-23.3$ & 0.202 & 2.11 & 3.54 & 7.13 & 2.90 & 0.99 & 0.16 & $0^\dagger$ & 3.17 & 1.77\\
\bottomrule
\end{tabular}
\end{table*}

\section{Discussion and limitations}
The results rest on a single end-of-day snapshot of delayed quotes; a panel of dates
would show how stable the regime contrast with 2011 is. Futures are inferred from
put--call parity rather than taken from traded VIX futures, although the parity fit
is tight. Nelder--Mead with random restarts is a local method: for the per-maturity fits,
different seeds reproduced the errors to two decimals, but the joint MRLRSV fit is not
converged. MRLRSVJ and the time-dependent $\sigma_t$ of the thesis are not
implemented. The loss is the thesis' MMLSE; an implied-volatility loss would weight the
wings differently.

\section{Conclusion}
The main empirical claims of Bao (2013) hold up in a very different market, fifteen
years later. Log-normal mean reversion alone cannot price VIX options, and upward
jumps or stochastic vol-of-vol, not both, are enough to fit a single maturity to
within the bid--ask spread. What the thesis left out matters as much. The fitted
parameters sit on flat ridges and bounds, the 2026 smile implies rare large jumps
rather than the frequent small ones of 2011, and no constant-parameter member of
the family explains the term structure. For pricing a given expiry, MRLRJ is a
compact and accurate smile model; as a model of how the VIX moves, the family needs
maturity-dependent parameters or richer dynamics.

\paragraph{Code and data.} {\raggedright The Julia package, tests, the 2026 snapshot and the
scripts that regenerate every table and figure are in \texttt{julia\_impl/}:
\texttt{Pkg.test()}, \texttt{scripts/reproduce\_paper.jl},
\texttt{scripts/calibrate\_2026.jl} and \texttt{scripts/make\_figures.jl}.\par}

{\footnotesize
\begin{thebibliography}{9}
\bibitem{bao2013} Q.~Bao. \emph{Mean-Reverting Logarithmic Modeling of VIX}. PhD
thesis, Zhejiang University, 2013. MPRA Paper No.~46413.
\bibitem{whaley1993} R.~Whaley. Derivatives on market volatility: hedging tools long
overdue. \emph{J.\ Derivatives}, 1:71--84, 1993.
\bibitem{grunbichler1996} A.~Gr\"unbichler, F.~Longstaff. Valuing futures and options on
volatility. \emph{J.\ Banking \& Finance}, 20:985--1001, 1996.
\bibitem{psychoyios2010} D.~Psychoyios, G.~Dotsis, R.~Markellos. A jump diffusion
model for VIX volatility options and futures. \emph{Rev.\ Quant.\ Finance Account.},
35(3):245--269, 2010.
\bibitem{gilpelaez1951} J.~Gil-Pelaez. Note on the inversion theorem.
\emph{Biometrika}, 38:481--482, 1951.
\bibitem{cboe} Cboe Global Markets. Delayed quotes, VIX index options.
\url{https://cdn.cboe.com/api/global/delayed_quotes/options/_VIX.json},
retrieved 29~Sep~2026.
\end{thebibliography}}
\end{document}
